\begin{figure}[H]
    \centering
    \begin{tikzpicture}[
        class/.style={draw, Navy, thick, minimum width=3cm, fill=white, rectangle split, rectangle split parts=2, font=\small},
        relation/.style={-Stealth, thick, Navy}
    ]
        \node[class] (league) {
            \textbf{League}
            \nodepart{second}
            + id: Int\\
            + name: String\\
            + country: String
        };

        \node[class, right=2cm of league] (season) {
            \textbf{Season}
            \nodepart{second}
            + year: Int\\
            + start\_date: Date\\
            + end\_date: Date
        };

        \node[class, below=2cm of season] (match) {
            \textbf{Match}
            \nodepart{second}
            + fixture\_id: Int\\
            + match\_date: DateTime\\
            + status: String\\
            + score\_home: Int\\
            + score\_away: Int
        };

        \node[class, left=2cm of match] (team) {
            \textbf{Team}
            \nodepart{second}
            + team\_id: Int\\
            + name: String\\
            + logo\_url: URL
        };

        \node[class, right=2cm of match] (stats) {
            \textbf{MatchStatistics}
            \nodepart{second}
            + possession: Float\\
            + shots: Int\\
            + passes\_acc: Float
        };

        \node[class, below=2cm of match] (lineup) {
            \textbf{MatchLineup}
            \nodepart{second}
            + home\_formation: Str\\
            + home\_xi: JSON\\
            + away\_xi: JSON
        };

        \draw[relation] (league) -- node[above, font=\tiny] {1..n} (season);
        \draw[relation] (season) -- node[right, font=\tiny] {1..n} (match);
        \draw[relation] (team) -- node[above, font=\tiny] {1..n (Home/Away)} (match);
        \draw[relation] (match) -- node[above, font=\tiny] {1..1} (stats);
        \draw[relation] (match) -- node[right, font=\tiny] {1..1} (lineup);

    \end{tikzpicture}
    \caption{Modèle Conceptuel de Données simplifié (ORM Django)}
\end{figure}
